\subsection{Several critical points with a common value}\label{sec:commondisks}

Choose the compact fibration of the preceding subsection at a $G$-fixed Morse parameter. Write $\pi$ for its map to the disk, $F_0$ for its marked part of the fibre over $R$, and $\lambda_1$ for the primitive after the horizontal boundary has been made a product. A thimble over an outgoing path is the Lagrangian disk swept out by a vanishing sphere under parallel transport. For a common critical value there is one such disk for each critical point, although all these disks project to the same path. We will need to order the disks in either direction while recognizing them as individual objects in a single symplectic manifold.

Choose one embedded path from each distinct critical value to an interval of small positive angles on the base circle, disjoint from $R$. Choose the paths pairwise disjoint, straight near their critical ends and radial near the base circle. At an admissible phase one can take parallel outgoing rays in the prescribed direction up to a large circle, then bend their disjoint tails to the interval in the same boundary order. Use the same path for every point over a common value.

\begin{lemma}\label{lem:samedisks}
There is a $G$-invariant exact change of primitive, supported away from $F_0$, for which the individual thimbles have cylindrical ends in the completion. For every refinement of the order at the common values, there is a compact exact Lefschetz fibration with distinct critical values on the same domain, having these very disks as thimbles in the prescribed order. Its total Liouville primitive, marked boundary fibre and completed disks can be kept fixed through all the refinements.
\end{lemma}

\begin{proof}
The local thimble consists of the points whose parallel transports towards the singular value converge to the specified critical point \cite[Lemma~1.13]{Sei}. At a common value take disjoint Morse neighbourhoods and a sufficiently short initial segment. The vanishing spheres at its regular end are disjoint in a single fibre. Transport along the rest of the path applies the same diffeomorphisms to these spheres. They stay disjoint and in the fibre interior, since the horizontal collar is a product. For different values disjointness follows from the paths. We therefore obtain pairwise disjoint compact disks $D_s$ meeting the vertical face away from the corners. Uniqueness of transport shows that $g D_s=D_{gs}$; compare \cite[\S13.1]{CH}. In complex dimension one the vanishing sphere consists of two points and their traces give the two ends of the same interval.

The disks are Lagrangian for $d(\lambda_1+\kappa\pi^*\alpha_0)$. Adding base area leaves the horizontal distribution unchanged, and its restriction to a disk over a path is zero. Hence $\lambda_1|_{D_s}=dh_s$ for a function on the disk. Extend these functions from terminal collars to disjoint tubular neighbourhoods of the ends up to the vertical face. Cut them off away from the horizontal face, the critical points and the marked fibre. Averaging the sum over $G$ gives a function $h$ with the same restricted differential on a smaller collar of every end. There $\pi^*\alpha_0|_{TD_s}=0$ by radiality, so
\[
               \Lambda=\lambda_1-dh+\kappa\pi^*\alpha_0,\qquad
               \Lambda|_{TD_s}=0\quad\text{on the terminal collars}.
\]
The Liouville field for $\Lambda$ is tangent to the disks on those collars because they are Lagrangian. On the support of $dh$ its horizontal contribution to the primitive is bounded independently of $\kappa$. By the vertical-face estimate of the preceding proof, for large $\kappa$ the projected Liouville field remains outward during subtraction of $t\,dh$, $0\le t\le1$. The disks thus meet the boundary along Legendrians and extend cylindrically, using the same rounding for all of them.

In a thin strip along each path $\gamma$ choose a transverse base vector field $V_\gamma$, constant in affine coordinates near its critical endpoint and equal to angular rotation near the outer circle. Choose compatible transverse signs at the ends and interpolate through positive transverse fields. Extend by cutoff to a field tangent to the circle with support away from the marked point. For sufficiently small times its translates of the whole path are disjoint by the tubular neighbourhood theorem. In the total space take disjoint tubes about the closed disks, including their critical points, and smooth cutoffs $\eta_s$ equal to one on neighbourhoods of the disks. For small distinct times $\epsilon_s$ within each group with common path, define
\[
             \pi'(p)=\operatorname{Fl}^{\,\epsilon_s\eta_s(p)}_{V_\gamma}
                                  \bigl(\pi(p)\bigr)
             \qquad (p\text{ in the tube of }D_s).
\]
Use $\pi$ elsewhere. Near each critical point this map is an affine translation of $\pi$. Away from smaller Morse balls it is arbitrarily close in $C^1$ to a symplectic submersion on a compact set, with respect to the fixed $d\Lambda$. Hence there are no new critical points and its regular fibres remain symplectic. The flow preserves the base circle. The horizontal product collar and a neighbourhood of the marked fibre are unchanged. Taking the $\epsilon_s$ in the desired transverse order gives the distinguished paths $\operatorname{Fl}_{V_\gamma}^{\epsilon_s}(\gamma)$.

For the fibrewise primitive use $\Lambda-(\pi')^*(\alpha_R+\kappa\alpha_0)$. This gives the product primitive on the horizontal collar. Near the critical points the translation preserves base area, so the local Kähler form from the preceding subsection is unchanged. On the neighbourhood where $\eta_s=1$, the differentials of $\pi$ and $\pi'$ have the same kernel. Their symplectic connections for $d\Lambda$ therefore have the same horizontal spaces there. The old trajectories on $D_s$ are precisely the new lifts of their translated paths, with the same collapse at the same point. This identifies the actual disks up to and including their ends. Any previously chosen brane gradings and relative Pin structures on them can be carried identically through the construction.
\end{proof}

This argument explains how directedness for separated thimbles can imply orthogonality at a repeated value. If the backward morphisms between the separated disks vanish in the partially wrapped category stopped at the marked fibre, order two specified members of a common fibre first one way and then the other. The two vanishings concern the same pair of objects in the same category by Lemma~\ref{lem:samedisks}. Applying this argument requires directedness in that category; disjoint compact vanishing spheres in one fibre alone compute a different Floer group.
